\section{Christian Grief, Burial of the Dead, and Prayer for the Departed}

Christian faith never commands the bereaved to turn their hearts to stone. Jesus wept before Lazarus's tomb (John 11:35); St.\ Ambrose wept over Satyrus; St.\ Augustine let his tears flow before God for Monica. What the Gospel forbids is not human sorrow, but hopelessness. St.\ Paul sets the measure for every Christian funeral:

\begin{quote}
And we will not have you ignorant brethren, concerning them that are asleep, that you be not sorrowful, even as others who have no hope. For if we believe that Jesus died and rose again: even so them who have slept through Jesus, will God bring with him. (1 Thess.\ 4:12--13)\endnote{1 Thess.\ 4:12--13 in Douay-Rheims (4:13--14 in the Clementine Vulgate; vv.\ 15--16 below correspond to Vulgate vv.\ 16--17).}
\end{quote}

St.\ John Chrysostom notices Paul's delicate first mercy: he ``has not said `concerning the dead,' but `them that are asleep,' even at the beginning suggesting consolation to them.''\endnote{St.\ John Chrysostom, \emph{Homilies on the First Epistle to the Thessalonians}, Hom.\ 6, on 1 Thess.\ 4:13 (NPNF, ser.\ 1, vol.\ 13).} Frantic pagan wailing---beating the breast and tearing the cheeks---belongs to ``a soul that knows nothing of the Resurrection, but thinks that this death is death.'' The Christian feels separation keenly, yet mourns inside the promise that the Lord will descend with the archangel's voice and God's trumpet, the dead in Christ will rise first, ``and so shall we be always with the Lord'' (1 Thess.\ 4:15--16). For Chrysostom, rightly ordered charity steadies the bereaved heart: ``If thou lovest God more, mourn not: for He who is more beloved is immortal, and does not suffer thee to feel the loss of him who is less beloved.''

\subsection{Reverence for the Dead Body}

Knowing what the body is and what it is destined to become, the Church treats the dead body neither with superstitious dread nor with callous disposal. St.\ Paul reminds the Corinthians that their members are ``the temple of God'' and ``the temple of the Holy Ghost,'' bought with a great price (1 Cor.\ 3:16--17; 6:19--20). In the Old Covenant, Tobias risked royal wrath to hide slain Israelites by day and bury them by night, and the Archangel Raphael declared that this work of mercy ascended before God with his prayers and tears (Tob.\ 1:20--21; 12:12). Hence the Church has always numbered burying the dead among the seven corporal works of mercy.

When Rome fell in 410 and pagans mocked Christians whose slain lay unburied, St.\ Augustine stated the exact Catholic balance in Book I of \emph{The City of God}. First, involuntary denial of burial never harms a righteous soul or hinders the Resurrection: ``there are indeed many bodies of Christians lying unburied; but no one has separated them from heaven, nor from that earth which is all filled with the presence of Him who knows whence He will raise again what He created.'' Pompous tombs profit the wicked dead nothing, whereas the beggar Lazarus was borne aloft by angels; funeral ceremonies ``are rather the solace of the living than the comfort of the dead.''\endnote{St.\ Augustine, \emph{The City of God}, I.12--13 (Dods, vol.\ 1); in I.13, the Latin text reads \emph{organis et uasis} (``organs and vessels''), CSEL 40.1.} Yet immediately Augustine guards the complementary truth:

\begin{quote}
Nevertheless the bodies of the dead are not on this account to be despised and left unburied; least of all the bodies of the righteous and faithful, which have been used by the Holy Ghost as His organs and instruments for all good works. For if the dress of a father, or his ring, or anything he wore, be precious to his children, in proportion to the love they bore him, with how much more reason ought we to care for the bodies of those we love, which they wore far more closely and intimately than any clothing! For the body is not an extraneous ornament or aid, but a part of man's very nature. (I.13)
\end{quote}

Recalling Tobias, the woman who anointed Christ's limbs for burial, and Joseph of Arimathea wrapping His Body in linen, Augustine concludes that ``God's providence extends even to the bodies of the dead, and that such pious offices are pleasing to Him, as cherishing faith in the resurrection.'' The Catholic cemetery takes its name from the Greek \emph{koimeterion}---a sleeping-place where mortal flesh rests in consecrated ground awaiting the harvest.

\subsection{The Requiem Liturgy and the Office of the Dead}

The Church's theology of sorrow and resurrection finds its supreme voice in her rites for the faithful departed. At the heart of the Requiem Mass, the priest lifts up the Preface for the Dead:

\begin{quote}
\emph{In quo nobis spes beat\ae{} resurrectionis effulsit, ut quos contristat certa moriendi conditio, eosdem consoletur futur\ae{} immortalitatis promissio. Tuis enim fidelibus, Domine, vita mutatur, non tollitur, et, dissoluta terrestris hujus incolatus domo, \ae{}terna in c\ae{}lis habitatio comparatur.}\endnote{\emph{Missale Romanum} (Tours: Alfred Mame, 1922), \emph{Praefatio pro Defunctis}, p.\ 339 (matching the 1962 typical edition).}
\end{quote}

``In whom the hope of a blessed resurrection hath shone upon us, that those whom the certain condition of dying saddens, the promise of future immortality may console. For unto Thy faithful, O Lord, life is changed, not taken away; and the house of this earthly sojourn being dissolved, an eternal dwelling is prepared in heaven.'' Natural grief (\emph{contristat certa moriendi conditio}) and supernatural hope (\emph{vita mutatur, non tollitur}) stand together without canceling one another.

The Office of the Dead echoes that chord through every Hour.\endnote{\emph{The Little Office of the Blessed Virgin Mary and the Office of the Dead} (New York: Benziger Brothers, 1915), Office of the Dead, Vespers, pp.\ 260--262, and Matins, pp.\ 285--297; cf.\ \emph{Manual of Prayers} (1889), p.\ 542.} Vespers opens with Psalm 114 and the antiphon \emph{Placebo Domino in regione vivorum} (``I will please the Lord in the land of the living''). Matins begins with the Invitatory \emph{Regem, cui omnia vivunt, venite adoremus} (``The King, unto whom all live: O come, let us worship Him'') and Psalm 5 under the antiphon \emph{Dirige, Domine, Deus meus, in conspectu tuo viam meam} (``Direct, O Lord, my God, my steps in Thy sight''), whence the English word ``dirge.'' Its nine Lessons from Job spare nothing of dust and darkness, yet after the very first Lesson the choir answers with Job 19:25--27: \emph{Credo quod Redemptor meus vivit; et in novissimo die de terra surrecturus sum. Et in carne mea videbo Deum Salvatorem meum}---``I believe my Redeemer liveth, and that in the last day I shall rise from the earth. And in my flesh I shall see God my Saviour.'' Throughout the Office, every Psalm closes not with the \emph{Gloria Patri} but with \emph{Requiem aeternam dona eis, Domine: et lux perpetua luceat eis}.

After the Requiem Mass, the priest stands before the bier for the Absolution, praying \emph{Non intres in judicium cum servo tuo, Domine}---``Enter not into judgment with Thy servant, O Lord\ldots{} who, while he lived, was sealed with the seal of the Holy Trinity''---sprinkling the coffin with holy water in memory of Baptism and incensing it in honor of the temple of the Holy Ghost.\endnote{\emph{Manual of Prayers} (1889), The Burial of the Dead, pp.\ 581--585.} As the body is borne out toward the grave, the choir chants \emph{In paradisum}:

\begin{quote}
\emph{In Paradisum deducant te Angeli, in tuo adventu suscipiant te Martyres, et perducant te in civitatem sanctam Jerusalem. Chorus Angelorum te suscipiat, et cum Lazaro quondam paupere aeternam habeas requiem.} (pp.\ 583--584)
\end{quote}

``May the Angels lead thee into Paradise; at thy coming may the Martyrs receive thee, and bring thee into the holy City, Jerusalem. May the Choir of Angels receive thee, and with Lazarus, once a beggar, mayest thou have eternal rest.'' At the grave, the priest asks God to station His holy angel as keeper of the tomb and chants Christ's promise at Bethany: \emph{Ego sum resurrectio et vita}---``I am the Resurrection and the Life: he that believeth in Me, although he be dead, shall live: and every one who liveth, and believeth in Me, shall never die'' (John 11:25--26).

\subsection{Prayer and Sacrifice for the Departed}

Christian mourning never halts at the cemetery gate. Because death does not sever the communion of saints, love for the departed immediately becomes intercession. Already Judas Machabeus sent twelve thousand drachms of silver to Jerusalem for sacrifice for his fallen soldiers, ``thinking well and religiously concerning the resurrection,'' and Scripture draws the explicit lesson: ``It is therefore a holy and wholesome thought to pray for the dead, that they may be loosed from sins'' (2 Mach.\ 12:43, 46).

From apostolic times, the Church carried that intercession directly to the Eucharistic altar. Instructing the newly baptized in fourth-century Jerusalem, St.\ Cyril explains that immediately after the Consecration the Church commemorates all who have fallen asleep, confident that offering Christ sacrificed for our sins propitiates our merciful God on their behalf as well as our own, far surpassing any crown subjects might weave to win a king's pardon for banished offenders.\endnote{St.\ Cyril of Jerusalem, \emph{Catechetical Lectures} (Mystagogical Catechesis V), 23.9--10 (NPNF, ser.\ 2, vol.\ 7).} St.\ John Chrysostom similarly bears witness that the Apostles themselves appointed the remembrance of the dead in the Eucharistic Mysteries, knowing the great benefit it wins for departed souls.\endnote{St.\ John Chrysostom, \emph{Homilies on Philippians}, Hom.\ 3 (NPNF, ser.\ 1, vol.\ 13).}

The Church solemnly defined this constant faith at the Council of Florence and the Council of Trent, teaching that the faithful departed still undergoing purification are truly assisted by the suffrages of the living---principally the propitiatory Sacrifice of the Mass, alongside prayers, almsgiving, and other works of piety.\endnote{Council of Florence, Bull of Union \emph{Laetentur caeli} (6 July 1439), DS 1304; Council of Trent, Session XXII (17 September 1562), Doctrine on the Sacrifice of the Mass, ch.\ 2 and canon 3 (DS 1743, 1753), and Session XXV (3--4 December 1563), \emph{Decree on Purgatory} (DS 1820). Full exposition of Purgatory, suffrages, and indulgences belongs to Part II of this series (\emph{The Books Were Opened}).} That is why St.\ Monica asked only to be remembered at the altar, and why every Roman Mass names the departed in the \emph{Memento etiam}. Prayer for the dead refuses to leave the living helpless beside an open grave, turning sorrow into active charity offered at the altar of the Risen Lord.
